\begin{abstract}
  In this work, we study various graph partitioning problems under a general demand model. In each such task, we are given a graph $G=(V,E,c,w)$ with a capacity function $c\colon E\to \mathbb{N}$ and a demand function $w\colon V\times V\to \mathbb{N}$. Our main focus is the problem of finding a cut $(S, \bar{S})$ minimizing the quantity
  \[
  \psi_w( S ) = \frac{c( S, \bar{S} )}{w( S, V )\cdot w( \bar{S}, V )}.
  \]
  Here, $c( S, \bar{S} )$ is the cost of edges between $S$ and the complement of $S$, $\bar{S}$, and $w( S, V )=w( S )+w( S, \bar{S} )$ is the sum of the internal demand within $S$, $w( S )$, and the demand between vertices of $S$ and $\bar{S}$, $w( S, \bar{S} )$. We call $\psi_w( S )$ the \emph{generalized conductance} of the cut $(S, \bar{S})$, and the task of minimizing $\psi_w( S )$ the \textsc{Generalized Conductance Problem}. Our main contribution is an algorithm with an $\mathcal{O}(\log n)$-approximation guarantee for this objective. Our result is achieved via a two-way reduction: first to the well-known \textsc{Generalized $k$-Multicut Problem}, and then to a constrained variant of the classic \textsc{Sparsest-Cut Problem}, with an additional upper-bound constraint on the amount of demand that may be cut.
  
  Moreover, we show that the above procedure can be leveraged to obtain an $\mathcal{O}(\log n)$-bicriteria approximation for \textsc{Graph Partitioning with Demands}, where the goal is to find a minimum-cost subset of edges $C$ such that for every component $H$ of $G\setminus C$, $w( H )\leq \rho\cdot w( V )$. This, in turn, yields an $\mathcal{O}(\log n)$-approximation for \textsc{Hierarchical Clustering with Demands}, the problem of finding a hierarchy of cuts that partitions the graph into increasingly refined clusters. For multiplicative demand functions, our framework improves these guarantees to $\mathcal{O}(\sqrt{\log n})$ and for trees we obtain an $\mathcal{O}(1)$-approximation for all of the above objectives.
\end{abstract}

\noindent\textbf{Keywords:} Graph partitioning with demands, Generalized conductance, Sparsest cut, Generalized multicut, Hierarchical clustering, Approximation algorithms.
